% !TeX program = pdflatex
\documentclass[aspectratio=169,xcolor={dvipsnames,table}]{beamer}

\input{../../_en_preambulo_beamer}
% Comandos y macros estadísticas compartidas para las presentaciones Beamer en inglés (Metropolis)

% Conjuntos numéricos básicos
\newcommand{\R}{\mathbb{R}}
\newcommand{\N}{\mathbb{N}}
\newcommand{\Q}{\mathbb{Q}}
\newcommand{\Z}{\mathbb{Z}}
\newcommand{\set}[1]{\left\{ #1 \right\}}
\newcommand{\abs}[1]{\left|#1\right|}
\newcommand{\norm}[1]{\left\| #1 \right\|}

% Comandos estadísticos y probabilísticos del curso
\newcommand{\Var}{\operatorname{Var}}
\newcommand{\cov}{\operatorname{Cov}}
\newcommand{\corr}{\rho}
\newcommand{\comb}[2]{\begin{pmatrix} #1 \\ #2 \end{pmatrix}}
\newcommand{\s}{\sigma}
\newcommand{\particion}{\mathfrak{P}}
\newcommand{\card}[1]{\#\left( #1 \right)}
\newcommand{\seq}[3]{\set{#1}_{#2}^{#3}}
\newcommand{\E}{\mathbb{E}}
\newcommand{\Prob}{\mathbb{P}}

% Entornos pedagógicos y cajas en inglés académico natural (soportando tanto nombres en inglés como en español)
\newenvironment{discussion}{%
  \begin{alertblock}{Discussion Question}%
}{%
  \end{alertblock}%
}
\newenvironment{discusion}{%
  \begin{alertblock}{Discussion Question}%
}{%
  \end{alertblock}%
}

\newenvironment{reflection}{%
  \begin{alertblock}{Classroom Reflection}%
}{%
  \end{alertblock}%
}
\newenvironment{reflexion}{%
  \begin{alertblock}{Classroom Reflection}%
}{%
  \end{alertblock}%
}

\newenvironment{paradox}{%
  \begin{alertblock}{Exploratory Paradox}%
}{%
  \end{alertblock}%
}
\newenvironment{paradoja}{%
  \begin{alertblock}{Exploratory Paradox}%
}{%
  \end{alertblock}%
}

\newenvironment{motivation}{%
  \begin{exampleblock}{Motivation: Data Science in Practice}%
}{%
  \end{exampleblock}%
}
\newenvironment{motivacion}{%
  \begin{exampleblock}{Motivation: Data Science in Practice}%
}{%
  \end{exampleblock}%
}

\newenvironment{quickchallenge}{%
  \begin{block}{Quick Team Challenge}%
}{%
  \end{block}%
}
\newenvironment{retoclase}{%
  \begin{block}{Quick Team Challenge}%
}{%
  \end{block}%
}


\title{Confidence Intervals for Population Means}
\subtitle{Section 06.04 --- Construction with Z and t, and Frequentist Interpretation}
\author[J. Castillo Colmenares]{Juliho Castillo Colmenares}
\institute[Tec de Monterrey]{Tecnologico de Monterrey}
\date{\vspace{-1.5cm}}

\begin{document}

% SLIDE 01: Title Page
\begin{frame}[plain]
	\titlepage
\end{frame}

% SLIDE 02: Roadmap
\begin{frame}{Roadmap --- Chapter 06: Estimation and Data Science}
	\vspace{-0.15cm}
	\begin{block}{Curricular Structure of Unit 5}
		\footnotesize
		\begin{enumerate}\itemsep=0.03cm
			\item \textcolor{gray}{Section 06.01: Point Estimation, Unbiasedness, Efficiency, and Consistency}
			\item \textcolor{gray}{Section 06.02: Method of Moments (MoM)}
			\item \textcolor{gray}{Section 06.03: Maximum Likelihood Estimation (MLE) and Score}
			\item \textbf{\textcolor{TecRojo}{Section 06.04: Confidence Intervals for Population Means ($Z$ and $t$)}} $\leftarrow$ \emph{Current focus}
			\item \textcolor{gray}{Section 06.05: Confidence Intervals for Variances ($\chi^2$) and Proportions}
		\end{enumerate}
	\end{block}
	\vspace{-0.2cm}
	\begin{alertblock}{Learning Objective}
		\scriptsize
		Apply the machinery of point estimation and sampling distributions to construct exact confidence intervals for $\mu$ --- using $Z$ when $\sigma$ is known, $t$ when it is not --- and rigorously understand what ``confidence'' means from the frequentist perspective.
	\end{alertblock}
\end{frame}

% SLIDE 03: Motivation
\begin{frame}{From Theory to Practice: A Single Number Is Not Enough}
	\footnotesize
	\begin{motivacion}
		A point estimator $\bar X$ gives a single number, without indicating how reliable it is. A confidence interval communicates \emph{both} things: a central estimate and a plausible range, calibrated by the exact sampling distribution of $\bar X$ (Chapter 05).
	\end{motivacion}
	\vspace{0.15cm}
	\pause
	\begin{itemize}\itemsep=0.1cm
		\item \textbf{Universal structure:} estimator $\pm$ (critical value) $\times$ (standard error). \pause
		\item \textbf{Known $\sigma$:} use $Z$ (rare in practice, but pedagogically fundamental). \pause
		\item \textbf{Unknown $\sigma$:} use $t$ with $n-1$ degrees of freedom (the common case).
	\end{itemize}
\end{frame}

% SLIDE 03B: Z vs t Intervals
\begin{frame}{Confidence Intervals for $\mu$: $Z$ vs. $t$}
	\vspace{-0.15cm}
	\begin{columns}[T]
		\begin{column}{0.48\textwidth}
			\begin{block}{Known $\sigma$}
				\footnotesize
				$$\bar X \pm Z_{\alpha/2}\dfrac{\sigma}{\sqrt n}$$
			\end{block}
		\end{column}
		\begin{column}{0.48\textwidth}
			\begin{alertblock}{Unknown $\sigma$ (Section 05.04)}
				\footnotesize
				$$\bar X \pm t_{\alpha/2,\,n-1}\dfrac{S}{\sqrt n}$$
				Always wider than the $Z$-interval.
			\end{alertblock}
		\end{column}
	\end{columns}
\end{frame}

% SLIDE 04: Common Structure and Extension
\begin{frame}{The Common Structure and Its Extension}
	\vspace{-0.15cm}
	\begin{block}{Estimator $\pm$ Margin of Error}
		\footnotesize
		Every confidence interval in this chapter shares the form: $\text{Point estimator} \pm (\text{critical value}) \times (\text{standard error})$.
	\end{block}
	\pause
	\vspace{0.1cm}
	\begin{alertblock}{Difference of Two Means}
		\footnotesize
		Known variances: $Z_{\alpha/2}\sqrt{\sigma_1^2/n_1+\sigma_2^2/n_2}$. Unknown equal variances: pooled $t$ with $S_p^2$, $\nu=n_1+n_2-2$. Unknown unequal variances: Welch-Satterthwaite. Paired samples: reduce to differences $D_i$.
	\end{alertblock}
\end{frame}

% SLIDE 05: Frequentist Interpretation
\begin{frame}{What Does ``Confidence'' Really Mean?}
	\vspace{-0.15cm}
	\begin{block}{Correct Interpretation (Ex-Ante)}
		\footnotesize
		$P(L(X) \le \mu \le U(X)) = 1-\alpha$ describes the \textbf{procedure}, before observing data: $\bar X$, $L(X)$, $U(X)$ are random variables.
	\end{block}
	\pause
	\vspace{0.1cm}
	\begin{alertblock}{The Common Mistake (Ex-Post)}
		\footnotesize
		Once a specific numeric interval like $[74.08, 81.92]$ is computed, $\mu$ either is or isn't in it: no probability is involved. ``95\% confidence'' describes the success rate of the \emph{method} under infinite repeated sampling.
	\end{alertblock}
\end{frame}

% SLIDE 06: Applications
\begin{frame}{Applications in Data Science and Engineering}
	\vspace{-0.1cm}
	\begin{itemize}\itemsep=0.1cm
		\item \textbf{Quality control:} CI for the average content of a production batch. \pause
		\item \textbf{A/B testing:} CI for the difference of means between experiment variants. \pause
		\item \textbf{Model reporting:} regression coefficient standard errors translate directly into CIs. \pause
		\item \textbf{Bayesian contrast:} the Bayesian credible interval does allow direct probabilistic statements about $\theta$, unlike the frequentist CI.
	\end{itemize}
\end{frame}

% SLIDE 07: Python Lab Block 1
\begin{frame}[fragile]{Python Lab: $Z$-Based vs. $t$-Based CI}
	\vspace{-0.05cm}
	\scriptsize
	Comparing intervals for a mean with known and unknown $\sigma$ (\texttt{06.04\_confidence\_intervals\_means.py}):
	{\lstset{basicstyle=\fontsize{4pt}{4.8pt}\selectfont\ttfamily,aboveskip=0.1em,belowskip=0.1em}
	\lstinputlisting[firstline=17, lastline=39]{../../code/06_estimacion_estadistica/06.04_confidence_intervals_means.py}}
\end{frame}

% SLIDE 08: Python Lab Block 2
\begin{frame}[fragile]{Python Lab: Pooled Two-Sample Mean-Difference CI}
	\vspace{-0.05cm}
	\scriptsize
	Computing the pooled variance $S_p^2$ and the CI for $\mu_A-\mu_B$:
	{\lstset{basicstyle=\fontsize{4pt}{4.8pt}\selectfont\ttfamily,aboveskip=0.1em,belowskip=0.1em}
	\lstinputlisting[firstline=42, lastline=61]{../../code/06_estimacion_estadistica/06.04_confidence_intervals_means.py}}
\end{frame}

% SLIDE 09: Python Lab Block 3
\begin{frame}[fragile]{Python Lab: Frequentist Coverage}
	\vspace{-0.05cm}
	\scriptsize
	Monte Carlo verification of the frequentist interpretation via repeated sampling:
	{\lstset{basicstyle=\fontsize{4pt}{4.8pt}\selectfont\ttfamily,aboveskip=0.1em,belowskip=0.1em}
	\lstinputlisting[firstline=64, lastline=89]{../../code/06_estimacion_estadistica/06.04_confidence_intervals_means.py}}
\end{frame}

% SLIDE 10: Python Lab Terminal Output
\begin{frame}[fragile]{Python Lab: Terminal Output}
	\vspace{-0.1cm}
	\begin{block}{}
		\fontsize{4pt}{5pt}\selectfont
		\begin{verbatim}
=== Block 1: Z-Based vs t-Based CI ===
Known sigma=2.0: 95% CI [499.42, 501.38]
Unknown sigma (S=2.3): 95% CI [499.17, 501.63] (25.1% wider)

=== Block 2: Pooled Two-Sample t CI ===
Sp^2=20.95, SE=1.96, df=20, diff=4.00
95% CI: [-0.088, 8.088] -- contains zero? Yes

=== Block 3: Frequentist Coverage ===
Empirical coverage: 0.9496 (nominal: 0.95)
		\end{verbatim}
	\end{block}
\end{frame}


% SLIDE 19: Closing and Chapter Perspective
\begin{frame}{Closing Section and Chapter Perspective}
	\vspace{-0.2cm}
	\begin{block}{Synthesis: Confidence Intervals for Means}
		\scriptsize
		Building a confidence interval is not just applying a formula: it requires choosing the correct reference distribution ($Z$ or $t$), understanding its common structure (estimator $\pm$ margin), and above all, correctly interpreting what ``confidence'' guarantees --- a property of the method, not of the particular interval obtained.
	\end{block}
	\vspace{0.05cm}
	\pause
	\begin{alertblock}{Key Lessons and Next Step}
		\begin{itemize}\itemsep=0.05cm
			\item \textbf{Known $\sigma$:} $\bar X \pm Z_{\alpha/2}\sigma/\sqrt n$. \textbf{Unknown $\sigma$:} $\bar X \pm t_{\alpha/2,n-1}S/\sqrt n$.
			\item \textbf{Mean difference:} pooled variance (equal) or Welch (unequal).
			\item \textbf{Interpretation:} confidence is a property of the procedure under repeated sampling, not a probability of the already-computed interval.
			\item \textbf{Next Section:} \textcolor{TecRojo}{Confidence Intervals for Variances ($\chi^2$) and Proportions}.
		\end{itemize}
	\end{alertblock}
\end{frame}

% SLIDE 20: Modular Perspective
\begin{frame}{Modular Perspective --- The Next Challenge}
  \small
  \begin{block}{What We Have Built}
    Confidence intervals quantify uncertainty for means and differences between means.
  \end{block}
  \pause
  \begin{alertblock}{Next Step}
    The next sections extend the same logic to proportions, variances, and sample size.
  \end{alertblock}
\end{frame}\end{document}
